% \iffalse meta-comment
%
% hit-tikz-demo.dtx 是示例文档那张染色体示意图用的 TikZ 宏
%
% \fi
%
% \section{示例插图用的 TikZ 宏}
%
% \changes{v3.2a}{0000/00/00}{示例插图用的 TikZ 宏与配套颜色从 \file{hit-float.dtx}
%   与 \file{hit-thesis.dtx} 抽出来单独成文件}
% \changes{v3.2a}{0000/00/00}{原 hithesis.sty 中的示例 tikz 样式与染色体示意图绘制宏并入 book-floats}
% \file{examples/hitbook/chinese/body/introduction.tex} 里有一张染色体示意图，
% 画它要一批 \pkg{tikz} 样式和 \cs{gseg}、\cs{rcell}、\cs{tcellone} 这些宏。
% 这批宏跟排版规范没有关系，画的是固定内容（\cs{rcell} 直接画 A~C~T~C 四个碱基），
% 手册里也没有写进去，本来不该待在文档类里。
%
% 先单独成文件，不动接口。真要清理的话是把它们做成示例目录下的一个宏包，
% 由示例文档自己 \cs{usepackage}，但那样 \cs{rcell} 这些就不再由类提供，
% 是接口变更，得另外定。\file{testfiles/03-user-macros.lvt} 钉着 \cs{rcell}
% 由类提供，动接口时那条断言要跟着改。
%
% 只有学位论文类装了这些，报告类没有。
%
% \begin{macro}{\gseg,\gsegr,\rcell,\ncell,\tcellone,\tcelltwo,\tcellthree,\tcellfour,\tcellfive}
% 四个颜色只给下面这批 \cs{tikzset} 用，跟着一起搬过来。
%    \begin{macrocode}
%<*thesis-tikz-demo>
\definecolor{colorzero}{rgb}{0,~0,~0}
\definecolor{colorone}{rgb}{1,~0,~0}
\definecolor{colortwo}{rgb}{0,~0,~1}
\definecolor{colorthree}{rgb}{0,~1,~0}
\tikzset{maternal/.style={colorone}}
\tikzset{paternal/.style={colortwo}}
\tikzset{variant/.style={colorthree!80!colorzero}}
\tikzset{reference/.style={colorzero}}
\tikzset{aallele/.style={colorzero,rotate=90}}
\tikzset{ballele/.style={colorthree!80!colorzero,rotate=90}}
\tikzset{refseg/.style={colorzero,draw=colorzero,~opacity=0.2}}
\tikzset{mseg/.style={colorone,draw=colorone,~opacity=0.2}}
\tikzset{pseg/.style={colortwo,draw=colortwo,~opacity=0.2}}
\tikzset{vseg/.style={colorthree!80!colorzero,draw=colorthree!80!colorzero,~opacity=0.6}}
\tikzset{bncell/.style={draw=colorzero,opacity=0.2,line~width=2pt,~rounded~corners=1pt}}
\tikzset{btcell/.style={draw=colorone,opacity=0.6,~line~width=2pt,~rounded~corners=1pt}}
\tikzset{tncell/.style={colorzero,opacity=0.9}}
\tikzset{ttcell/.style={colorone,opacity=0.6}}
\tikzset{tscell/.style={colorzero}}
\tikzset{refcell/.style={colorzero}}
\tikzset{evolve/.style={->,draw=colortwo,opacity=0.3,line~width=1.5pt}}
\tikzset{fakeevolve/.style={->,draw=colorzero,opacity=0.3,line~width=1.5pt}}
\tikzset{refline/.style={dashed,draw=colorzero,line~width=1pt}}
\tikzset{tnline/.style={dashed,draw=colorzero,opacity=0.3,line~width=1pt}}
\NewDocumentCommand \gseg { m m m m m m m m m }
  {
%
~\pgfmathsetmacro{\sstartx}{#1}~
 \pgfmathsetmacro{\slengx}{#2}~
 \pgfmathsetmacro{\sy}{#3}~
 \pgfmathsetmacro{\sdy}{#4}~
 \pgfmathsetmacro{\sdx}{#5}~
 \pgfmathsetmacro{\sdxh}{#7}~
 \pgfmathsetmacro{\sdxt}{#8}~
 \fill[#6](\sstartx,\sy)--(\sstartx-\sdx,\sy+\sdy)--~
 (\slengx+\sstartx+1.5-\sdx,\sy+\sdy)--(\slengx+\sstartx+1.5,\sy)--~
 (\slengx+\sstartx+1.5-\sdx,\sy-\sdy)--(\sstartx-\sdx,\sy-\sdy)--cycle;~
 \draw[#9]~(\sstartx-\sdxh,\sy)~--~(\sstartx,~\sy);~
 \draw[#9]~(\slengx+\sstartx+1.5,~\sy)~--~(\slengx+\sstartx+1.5+\sdxt,\sy);~
  }
\NewDocumentCommand \gsegr { m m m m m m m m m }
  {
%
~\pgfmathsetmacro{\sstartx}{#1}~
 \pgfmathsetmacro{\slengx}{#2}~
 \pgfmathsetmacro{\sy}{#3}~
 \pgfmathsetmacro{\sdy}{#4}~
 \pgfmathsetmacro{\sdx}{#5}~
 \pgfmathsetmacro{\sdxh}{#7}~
 \pgfmathsetmacro{\sdxt}{#8}~
 \fill[#6](\sstartx-0.5,\sy)--(\sstartx+\sdx-0.5,\sy+\sdy)--~
 (\slengx+\sstartx+1.5+\sdx-0.5,\sy+\sdy)--(\slengx+\sstartx+1.5-0.5,\sy)--~
 (\slengx+\sstartx+1.5+\sdx-0.5,\sy-\sdy)--(\sstartx+\sdx-0.5,\sy-\sdy)--cycle;~
 \draw[#9]~(\sstartx-\sdxh-0.5,\sy)~--~(\sstartx-0.5,~\sy);~
 \draw[#9]~(\slengx+\sstartx+1.5-0.5,~\sy)~--~(\slengx+\sstartx+1.5+\sdxt-0.5,\sy);~
  }
\NewDocumentCommand \rcell { m m }
  {
%
~\pgfmathsetmacro{\x}{#1}~
 \pgfmathsetmacro{\y}{#2}~
 \draw (\x+1,\y)~node[aallele]{A};~
 \draw (\x+2,\y)~node[aallele]{C};~
 \draw (\x+3,\y)~node[aallele]{T};~
 \draw (\x+4,\y)~node[aallele]{C};~
 \gseg{\x}{4}{\y}{0.2}{0.5}{refseg}{1.5}{1}{reference}~
  }
\NewDocumentCommand \ncell { m m }
  {
%
~\pgfmathsetmacro{\x}{#1}~
 \pgfmathsetmacro{\y}{#2}~
 \draw[bncell](\x-2,\y+0.5)--(\x+7,\y+0.5)--~
 (\x+7,\y-1)--(\x-2,\y-1)--cycle;~
 \draw (\x+1,\y)~node[aallele]{A};~
 \draw (\x+2,\y)~node[ballele]{G};~
 \draw (\x+3,\y)~node[aallele]{T};~
 \draw (\x+4,\y)~node[aallele]{C};~
 \gseg{\x}{4}{\y}{0.2}{0.5}{mseg}{1.5}{1}{maternal}~
 \draw (\x+1,\y-0.5)~node[ballele]{T};~
 \draw (\x+2,\y-0.5)~node[aallele]{C};~
 \draw (\x+3,\y-0.5)~node[aallele]{T};~
 \draw (\x+4,\y-0.5)~node[ballele]{A};~
 \gseg{\x}{4}{\y-0.5}{0.2}{0.5}{pseg}{1.5}{1}{paternal}~
  }
\NewDocumentCommand \tcellone { m m }
  {
%
~\pgfmathsetmacro{\x}{#1}~
 \pgfmathsetmacro{\y}{#2}~
 \draw[btcell](\x-2,\y+0.5)--(\x+7,\y+0.5)--~
 (\x+7,\y-1.5)--(\x-2,\y-1.5)--cycle;~
 \draw (\x+1,\y)~node[aallele]{A};~
 \draw (\x+2,\y)~node[ballele]{G};~
 \draw (\x+3,\y)~node[aallele]{T};~
 \draw (\x+4,\y)~node[aallele]{C};~
 \gseg{\x}{4}{\y}{0.2}{0.5}{mseg}{1.5}{1}{maternal}~
 \draw (\x+1,\y-0.5)~node[aallele]{A};~
 \draw (\x+2,\y-0.5)~node[ballele]{G};~
 \draw (\x+3,\y-0.5)~node[aallele]{T};~
 \draw (\x+4,\y-0.5)~node[aallele]{C};~
 \gseg{\x}{4}{\y-0.5}{0.2}{0.5}{mseg}{1.5}{1}{maternal}~
 \draw (\x+1,\y-1)~node[ballele]{T};~
 \draw (\x+2,\y-1)~node[aallele]{C};~
 \draw (\x+3,\y-1)~node[aallele]{T};~
 \draw (\x+4,\y-1)~node[ballele]{A};~
 \gseg{\x}{4}{\y-1}{0.2}{0.5}{pseg}{1.5}{1}{paternal}~
  }
\NewDocumentCommand \tcellthree { m m }
  {
%
~\pgfmathsetmacro{\x}{#1}~
 \pgfmathsetmacro{\y}{#2}~
 \draw[btcell](\x-2,\y+0.5)--(\x+11,\y+0.5)--~
 (\x+11,\y-1)--(\x-2,\y-1)--cycle;~
 \draw (\x+1,\y)~node[aallele]{A};~
 \draw (\x+2,\y)~node[ballele]{G};~
 \gseg{\x}{2}{\y}{0.2}{0.5}{mseg}{1.5}{0}{maternal}~
 \gseg{\x+4}{0}{\y}{0.2}{0.5}{vseg}{0.5}{0.5}{variant}~
 \draw (\x+7,\y)~node[aallele]{T};~
 \draw (\x+8,\y)~node[aallele]{C};~
 \gseg{\x+6}{2}{\y}{0.2}{0.5}{mseg}{0}{1}{maternal}~
 \draw (\x+1,\y-0.5)~node[ballele]{T};~
 \draw (\x+2,\y-0.5)~node[aallele]{C};~
 \draw (\x+3,\y-0.5)~node[aallele]{T};~
 \draw (\x+4,\y-0.5)~node[ballele]{A};~
 \gseg{\x}{4}{\y-0.5}{0.2}{0.5}{pseg}{1.5}{1}{paternal}~
  }
\NewDocumentCommand \tcellfour { m m }
  {
%
~\pgfmathsetmacro{\x}{#1}~
 \pgfmathsetmacro{\y}{#2}~
 \draw[btcell](\x-2,\y+0.5)--(\x+15,\y+0.5)--~
 (\x+15,\y-1)--(\x-2,\y-1)--cycle;~
 \draw (\x+1,\y)~node[aallele]{A};~
 \draw (\x+2,\y)~node[ballele]{G};~
 \gseg{\x}{2}{\y}{0.2}{0.5}{mseg}{1.5}{0}{maternal}~
 \gseg{\x+4}{0}{\y}{0.2}{0.5}{vseg}{0.5}{0.5}{variant}~
 \draw (\x+7,\y)~node[aallele]{T};~
 \gseg{\x+6}{1}{\y}{0.2}{0.5}{mseg}{0}{0}{maternal}~
 \gseg{\x+9}{0}{\y}{0.2}{0.5}{vseg}{0.5}{0.5}{variant}~
 \draw (\x+12,\y)~node[aallele]{C};~
 \gseg{\x+11}{1}{\y}{0.2}{0.5}{mseg}{0}{1}{maternal}~
 \draw (\x+1,\y-0.5)~node[ballele]{T};~
 \draw (\x+2,\y-0.5)~node[aallele]{C};~
 \draw (\x+3,\y-0.5)~node[aallele]{T};~
 \draw (\x+4,\y-0.5)~node[ballele]{A};~
 \gseg{\x}{4}{\y-0.5}{0.2}{0.5}{pseg}{1.5}{1}{paternal}~
  }
\NewDocumentCommand \tcelltwo { m m }
  {
%
~\pgfmathsetmacro{\x}{#1}~
 \pgfmathsetmacro{\y}{#2}~
 \draw[btcell](\x-2,\y+0.5)--(\x+7,\y+0.5)--~
 (\x+7,\y-1.5)--(\x-2,\y-1.5)--cycle;~
 \draw (\x+1,\y)~node[aallele]{A};~
 \draw (\x+2,\y)~node[ballele]{G};~
 \draw (\x+3,\y)~node[aallele]{T};~
 \draw (\x+4,\y)~node[aallele]{C};~
 \gseg{\x}{4}{\y}{0.2}{0.5}{mseg}{1.5}{1}{maternal}~
 \draw (\x+1,\y-0.5)~node[aallele]{A};~
 \draw (\x+2,\y-0.5)~node[ballele]{G};~
 \draw (\x+3,\y-0.5)~node[aallele]{T};~
 \draw (\x+4,\y-0.5)~node[ballele]{G};~
 \gseg{\x}{4}{\y-0.5}{0.2}{0.5}{mseg}{1.5}{1}{maternal}~
 \draw (\x+1,\y-1)~node[ballele]{T};~
 \draw (\x+2,\y-1)~node[aallele]{C};~
 \draw (\x+3,\y-1)~node[aallele]{T};~
 \draw (\x+4,\y-1)~node[ballele]{A};~
 \gseg{\x}{4}{\y-1}{0.2}{0.5}{pseg}{1.5}{1}{paternal}~
  }
\NewDocumentCommand \tcellfive { m m }
  {
%
~\pgfmathsetmacro{\x}{#1}~
 \pgfmathsetmacro{\y}{#2}~
 \draw[btcell](\x-2,\y+0.5)--(\x+9.5,\y+0.5)--~
 (\x+9.5,\y-1.5)--(\x-2,\y-1.5)--cycle;~
 \draw (\x+1,\y)~node[aallele]{A};~
 \draw (\x+2,\y)~node[ballele]{G};~
 \draw (\x+3,\y)~node[aallele]{T};~
 \draw (\x+4,\y)~node[aallele]{C};~
 \gseg{\x}{4}{\y}{0.2}{0.5}{mseg}{1.5}{1}{maternal}~
 \draw (\x+1,\y-0.5)~node[aallele]{A};~
 \draw (\x+2,\y-0.5)~node[ballele]{G};~
 \draw (\x+3,\y-0.5)~node[aallele]{T};~
 \draw (\x+4,\y-0.5)~node[aallele]{C};~
 \gseg{\x}{4}{\y-0.5}{0.2}{0.5}{mseg}{1.5}{1}{maternal}~
 \draw (\x+1,\y-1)~node[ballele]{T};~
 \gseg{\x}{1}{\y-1}{0.2}{0.5}{pseg}{1.5}{0}{paternal}~
 \draw (\x+4.5,\y-1)~node[ballele]{A};~
 \draw (\x+5.5,\y-1)~node[aallele]{T};~
 \draw (\x+6.5,\y-1)~node[aallele]{C};~
 \gsegr{\x+3.5}{3}{\y-1}{0.2}{0.5}{pseg}{0.5}{1.5}{paternal}~
  }
%</thesis-tikz-demo>
%    \end{macrocode}
% \end{macro}
%
% \Finale
